% part: intuitionistic-logic
% chapter: soundness-completeness
% section: soundness-nd

\documentclass[../../../include/open-logic-section]{subfiles}

\begin{document}

\olfileid[pt]{int}{sc}{snd}

\olsection{Correção da dedução natural}

Provaremos agora a correção da dedução natural em relação à semântica
relacional, isto é, mostraremos que, se !!a{formula} é !!{derivable} a
partir de um conjunto de hipóteses, então esse conjunto de hipóteses
implica a !!{formula}.

\begin{thm}[Correção]
  \ollabel{thm:soundness}
  Se $\Gamma \Proves !A$, então $ \Gamma \Entails !A$.
\end{thm}

\begin{proof}
  Provamos que, se $\Gamma \Proves !A$, então $\Gamma \Entails !A$.
  A prova é por indução sobre a !!{derivation} de~$!A$ a partir
  de~$\Gamma$.
  
  \begin{enumerate}
  \item Se a !!{derivation} consiste apenas na hipótese~$!A$, temos
    $!A \Proves !A$ e queremos mostrar que $!A \Entails !A$. Suponha
    $\mSat{M}{!A}[w]$. Então, trivialmente, $\mSat{M}{!A}[w]$.

  \item A !!{derivation} termina em $\Intro{\land}$:
    \iftag{probAnd}{Exercício.}{As !!{derivation}s das premissas
      $!B$ a partir das hipóteses~$\Gamma$ !!{undischarged} e de~$!C$
      a partir das hipóteses~$\Delta$ !!{undischarged} mostram que
      $\Gamma \Proves !B$ e $\Delta \Proves !C$. Pela hipótese de
      indução, temos $\Gamma \Entails !B$ e $\Delta \Entails !C$.
      Precisamos mostrar que $\Gamma \cup \Delta \Entails !A \land !B$,
      pois as hipóteses !!{undischarged} de toda a !!{derivation} são
      $\Gamma$ juntamente com $\Delta$. Suponha, então,
      $\mSat{M}{\Gamma \cup \Delta}[w]$. Logo, também
      $\mSat{M}{\Gamma}[w]$. Como $\Gamma \Entails !B$,
      $\mSat{M}{!B}[w]$. De modo semelhante, $\mSat{M}{!C}[w]$.
      Portanto, $\mSat{M}{!B \land !C}[w]$.}

  \item A !!{derivation} termina em $\Elim{\land}$:
    \iftag{probAnd}{Exercício.}{A !!{derivation} da premissa
        $!B \land !C$ a partir das hipóteses $\Gamma$ !!{undischarged}
        mostra que $\Gamma \Proves !B \land !C$. Pela hipótese de
        indução, $\Gamma \Entails !B \land !C$. Precisamos mostrar que
        $\Gamma \Entails !B$. Suponha, então, $\mSat{M}{\Gamma}[w]$.
        Como $\Gamma \Entails !B \land !C$,
        $\mSat{M}{!B \land !C}[w]$. Logo, também $\mSat{M}{!B}[w]$.
        De modo semelhante, se $\Elim\land$ termina em~$!C$, então
        $\Gamma \Entails !C$.}

  \item A !!{derivation} termina em $\Intro{\lor}$:
    \iftag{probOr}{Exercício.}{Suponha que a premissa seja $!B$ e que
      as hipóteses !!{undischarged} da !!{derivation} que termina
      em~$!B$ sejam $\Gamma$. Então temos $\Gamma \Proves !B$ e, pela
      hipótese de indução, $\Gamma \Entails !B$. Precisamos mostrar que
      $\Gamma \Entails !B \lor !C$. Suponha $\mSat{M}{\Gamma}[w]$.
      Como $\Gamma \Entails !B$, $\mSat{M}{!B}[w]$. Mas então também
      $\mSat{M}{!B \lor !C}[w]$. De modo semelhante, se a premissa é
      $!C$, temos $\Gamma \Entails !C$.}
    
  \item A !!{derivation} termina em~$\Elim{\lor}$:
    \iftag{probOr}{Exercício.}{As !!{derivation}s que terminam nas
      premissas são de $!B \lor !C$ a partir das hipóteses~$\Gamma$
      !!{undischarged}, de $!D$ a partir das hipóteses
      $\Delta_1 \cup \{!B\}$ !!{undischarged} e de $!D$ a partir das
      hipóteses $\Delta_2 \cup \{!C\}$ !!{undischarged}. Portanto,
      temos $\Gamma \Proves !B \lor !C$,
      $\Delta_1 \cup \{!B\} \Proves !D$ e
      $\Delta_2 \cup \{!C\} \Proves !D$. Pela hipótese de indução,
      $\Gamma \Entails !B \lor !C$,
      $\Delta_1 \cup \{!B\} \Entails !D$ e
      $\Delta_2 \cup \{!C\} \Entails !D$. Precisamos provar que
      $\Gamma \cup \Delta_1 \cup \Delta_2 \Entails !D$.

    Suponha $\mSat{M}{\Gamma \cup \Delta_1 \cup \Delta_2}[w]$.
    Então $\mSat{M}{\Gamma}[w]$ e, como
    $\Gamma \Entails !B \lor !C$, $\mSat{M}{!B \lor !C}[w]$. Pela
    definição de $\mSat{M}{}$, ou $\mSat{M}{!B}[w]$ ou
    $\mSat{M}{!C}[w]$. Distinguimos, portanto, os casos:
    (a)~$\mSat{M}{!B}$[w]. Então
    $\mSat{M}{\Delta_1 \cup \{!B\}}[w]$. Como
    $\Delta_1 \cup !B \Entails !D$, temos $\mSat{M}{!D}[w]$.
    (b)~$\mSat{M}{!C}[w]$. Então
    $\mSat{M}{\Delta_2 \cup \{!C\}}[w]$. Como
    $\Delta_2 \cup !C \Entails !D$, temos $\mSat{M}{!D}[w]$. Assim,
    em ambos os casos, $\mSat{M}{!D}[w]$, como queríamos mostrar.}

  \item A !!{derivation} termina em $\Intro{\lif}$ com a conclusão
    $!B\lif !C$. A premissa é~$!C$, e a !!{derivation} que termina na
    premissa tem as hipóteses~$\Gamma \cup \{!B\}$ !!{undischarged}.
    Portanto, $\Gamma \cup \{!B\} \Proves !C$ e, pela hipótese de
    indução, $\Gamma \cup \{!B\} \Entails !C$. Precisamos mostrar que
    $\Gamma \Entails !B \lif !C$.

    Suponha $\mSat{M}{\Gamma}[w]$. Queremos mostrar que, para todo
    $w'$ tal que $Rww'$, se $\mSat{M}{!B}[w']$, então
    $\mSat{M}{!C}[w']$. Suponha, portanto, $Rww'$ e
    $\mSat{M}{!B}[w']$. Por \olref[sem][rel]{prop:true-monotonic},
    $\mSat{M}{\Gamma}[w']$. Como $\Gamma \cup \{!B\} \Entails !C$,
    $\mSat{M}{!C}[w']$, que é o que queríamos mostrar. 

  \item A !!{derivation} termina em $\Elim{\lif}$ e tem a
    conclusão~$!C$. As premissas são $!B \lif !C$ e $!B$, com
    !!{derivation}s a partir das hipóteses $\Gamma$, $\Delta$
    !!{undischarged}. Portanto, temos $\Gamma \Proves !B \lif !C$ e
    $\Delta \Proves !B$. Pela hipótese de indução,
    $\Gamma \Entails !B \lif !C$ e $\Delta \Entails !B$. Precisamos
    mostrar que $\Gamma \cup \Delta \Entails !C$.

    Suponha $\mSat{M}{\Gamma \cup \Delta}[w]$. Como
    $\mSat{M}{\Gamma}[w]$ e $\Gamma \Entails !B \lif !C$,
    $\mSat{M}{!B \lif !C}[w]$. Por definição, isso significa que, para
    todo $w'$ tal que $Rww'$, se $\mSat{M}{!B}[w']$, então
    $\mSat{M}{!C}[w']$. Como $R$ é reflexiva, $w$ está entre os $w'$
    tais que $Rww'$, O.s.\@, se $\mSat{M}{!B}[w]$, então
    $\mSat{M}{!C}[w]$. Como $\mSat{M}{\Delta}[w]$ e
    $\Delta \Entails !B$, $\mSat{M}{!B}[w]$. Logo,
    $\mSat{M}{!C}[w]$, como queríamos mostrar.

  \item A !!{derivation} termina em $\FalseInt$, concluindo~$!A$. A
    premissa é $\lfalse$, e as hipóteses !!{undischarged} da
    !!{derivation} da premissa são~$\Gamma$. Então
    $\Gamma \Proves \lfalse$.
    Pela hipótese de indução, $\Gamma \Entails \lfalse$. Precisamos
    mostrar $\Gamma \Entails !A$.

    Procedemos indiretamente. Se $\Gamma \Entails/ !A$, existe um
    modelo~$\mModel{M}$ e um mundo~$w$ tais que
    $\mSat{M}{\Gamma}[w]$ e $\mSat/{M}{!A}[w]$. Como
    $\Gamma \Entails \lfalse$, $\mSat{M}{\lfalse}[w]$. Mas isso é
    impossível, pois, por definição, $\mSat/{M}{\lfalse}[w]$.
    Portanto, $\Gamma \Entails !A$.
  \item A !!{derivation} termina em $\Intro\lnot$: exercício.
  \item A !!{derivation} termina em $\Elim\lnot$: exercício.
  \end{enumerate}
\end{proof}

\begin{prob}
  Complete a prova de \olref[int][sc][snd]{thm:soundness}. Nos casos
  de $\Intro\lnot$ e $\Elim\lnot$, use a definição de
  $\mSat{M}{\lnot !A}[w]$ em \olref[int][sem][rel]{defn:true-at-w}, O.s.\@, não trate $\lnot !A$ como definida por $!A \lif \lfalse$.
\end{prob}

\begin{prob}
  Mostre que as seguintes !!{formula}s não são !!{derivable} na lógica
  intuicionista:
  \begin{enumerate}
    \item $(!A \lif !B) \lor (!B \lif !A)$
    \item $(\lnot\lnot !A \lif !A) \lif (!A \lor \lnot !A)$
    \item $(!A \lif !B \lor !C) \lif \bigl((!A \lif !B)\lor(!A \lif !C)\bigr)$
  \end{enumerate}
\end{prob}

\end{document}
